% Pista geo-26s-q4b · Geometria
% Procedència: PAU Catalunya 2026 · Convocatòria de setembre · Sèrie 2 · Exercici 4 — Opció B
\documentclass[11pt,a4paper]{article}
\usepackage{pista}
\pistainfo{Catalunya 2026}{Convocatòria de setembre}{Sèrie 2}

\begin{document}

\section*{\textcolor{blauPista}{Exercici 4 --- Opció B}}

\noindent L'Aina vol construir una capsa de cartró amb forma de prisma rectangular amb tapa, per guardar-hi la tauleta i el carregador.

\noindent L'Aina vol fer un croquis en 3D de la capsa utilitzant un programa informàtic. Per a començar, situa un costat de la tapa des del vèrtex $A(1, 1, 1)$ fins al vèrtex $B(4, 1, 5)$, i l'altre costat de la tapa des del vèrtex $A$ fins al vèrtex $C(1, 10, 1)$.

\subsection*{4.1. \normalfont Quines són les coordenades del quart vèrtex de la tapa? Calculeu les longituds dels dos costats de la tapa. \hspace*{\fill}\textnormal{[1~punt]}}

\pista{Els dos costats de la tapa surten del vèrtex $A$: un va fins a $B$, i l'altre fins a $C$. Per tant, el quart vèrtex $D$ és el que queda davant del vèrtex $A$. Fes un dibuix del rectangle.

\medskip
Per arribar a $D$, has de partir de $B$ i aplicar-hi el mateix desplaçament que porta del vèrtex $A$ al vèrtex $C$:
\[
D = B + \vec{AC}
\]
Finalment, les longituds dels dos costats són, respectivament, els mòduls dels vectors $\vec{AB}$ i $\vec{AC}$.}

\subsection*{4.2. \normalfont Sabent que els altres quatre vèrtexs de la capsa estan situats al pla $\pi: -4x + 3z = 24$, calculeu les coordenades d'un d'aquests vèrtexs (qualsevol) i l'alçària de la capsa. \hspace*{\fill}\textnormal{[1,5~punts]}}

\pista{Aquest apartat et demana dues coses.

\medskip
\textit{Un vèrtex.} La capsa és un prisma recte: cada vèrtex de la base és just a sota d'un vèrtex de la tapa. Per trobar el vèrtex que és a sota del vèrtex $A$, escriu la recta que passa per $A$ amb el vector normal del pla $\pi$ com a vector director. Fixa't que a l'equació de $\pi$ no hi ha $y$: la 2a coordenada del vector normal és $0$.

\medskip
A continuació, troba el punt on aquesta recta talla el pla $\pi$.

\medskip
\textit{L'altura.} És la distància d'un vèrtex de la tapa, per exemple $A$, al pla $\pi$. En primer lloc, escriu el pla en la forma \mbox{$\ldots = 0$}, passant el $24$ a l'altre costat.

\medskip
A continuació, aplica la fórmula de la distància d'un punt $(x_{0}, y_{0}, z_{0})$ a un pla:
\[
d(A, \pi) = \frac{|a x_{0} + b y_{0} + c z_{0} + d|}{\sqrt{a^{2} + b^{2} + c^{2}}}
\]}

\end{document}
